\section{Threat Model}
\label{sec:threat}

\subsection{Adversary Model}

CPSForge models a \emph{post-compromise adversary} with S7 protocol-level network access to the PLC---achieved, e.g., through credential theft, insider access, or lateral movement from the IT network. This is realistic: S7 communication is commonly exposed on flat industrial networks, and tools such as python-snap7 allow any network-reachable host to read and write PLC data blocks without authentication on many deployed configurations.

\textbf{Capabilities.} The adversary can: (i)~continuously read the live tag stream (sensor values, actuator states, setpoints) at polling rate; (ii)~infer the current operational phase from tag patterns; (iii)~inject writes to writable tags in the PLC data block. Attacks are expressed at the \emph{process-semantic level}---actuator overrides, setpoint shifts, sensor spoofing, timing delays, and sequence perturbations---not firmware or memory-level modifications. The exact attack surface is scene-specific (Table~\ref{tab:attack_surface}).

\textbf{Goals.} Cause measurable physical impact (unsafe state transitions, degraded performance, incorrect sequencing) while remaining process-valid. When possible, maintain plausible sensor/actuator patterns to delay detection.

\subsection{LLM Agent Role}

The LLM operates as a man-in-the-middle in the tag stream, not as a trusted controller. Every three polling cycles (1.5\,s), it receives a structured snapshot of the observed tags and returns a structured decision: as red team, \emph{attack} or \emph{wait} (with a complete write specification); as blue team, \emph{allow} or \emph{block} (with suspicion score). All outputs are treated as untrusted proposals. The safety shield enforces tag whitelists, value ranges, duration caps, and process invariants before any write reaches the PLC.

\textbf{Timing.} LLM inference takes 13--33\,s per call (3B model, depending on context tier). This positions the adversary as a \emph{supervisory-timescale} attacker making strategic interventions at multi-second intervals---consistent with documented ICS incidents, where Stuxnet operated over weeks~\cite{langner2011stuxnet}, Ukraine grid attacks over minutes~\cite{liang2017ukraine}, and even HARVEY~\cite{garcia2017harvey} operated at 100\,ms--1\,s. The 10-attack budget tests whether the LLM can \emph{select the right writes} from limited opportunities.

\subsection{Attack Surface}
\label{sec:threat:attack_surface}

Table~\ref{tab:attack_surface} enumerates the concrete variables accessible via S7 protocol across our three evaluation scenes. All reside in non-optimized PLC data blocks readable/writable by python-snap7 without authentication. The ``PLC refresh'' column indicates whether the PLC overwrites the variable each scan cycle ($\checkmark$) or trusts the stored value ($\times$)---a critical distinction for attack persistence.

\begin{table}[t]
\centering
\caption{LLM-accessible attack surface per scene.}
\label{tab:attack_surface}
\small
\begin{tabular}{@{}llllc@{}}
\toprule
\textbf{Scene} & \textbf{Variable} & \textbf{Type} & \textbf{Category} & \textbf{Refresh} \\
\midrule
\multirow{4}{*}{\shortstack[l]{Level\\Control}} & SetpointIn & REAL & Setpoint hijack & $\times$ \\
& FillValve & REAL & Actuator override & $\checkmark$ \\
& DischargeValve & REAL & Actuator override & $\checkmark$ \\
& Enable & BOOL & Safety kill & $\times$ \\
\midrule
\multirow{5}{*}{\shortstack[l]{Sorting\\Weight}} & LightThresh & REAL & Threshold manip. & $\times$ \\
& HeavyThresh & REAL & Threshold manip. & $\times$ \\
& iState & INT & State jump & $\checkmark$ \\
& Conveyor1/2/3 & BOOL & Actuator override & $\checkmark$ \\
& Sorter1/2 & BOOL & Actuator override & $\checkmark$ \\
\midrule
\multirow{5}{*}{\shortstack[l]{Sorting\\Height}} & iState & INT & State jump & $\checkmark$ \\
& TurnTimer & INT & Timer accel. & $\checkmark$ \\
& DwellTimer & INT & Timer accel. & $\checkmark$ \\
& Enable & BOOL & Safety kill & $\times$ \\
& IsTall & BOOL & Flag spoof & $\checkmark$ \\
\bottomrule
\end{tabular}
\end{table}

\subsection{Defender Model}

The defender observes the same tag stream plus experiment telemetry. Depending on configuration, it uses static thresholds, process invariants, phase-aware rules, intent-consistency checks, or LLM-based anomaly detection (\S\ref{sec:design}). The defender has no oracle knowledge of attacker intent; logged attacks can be used for offline adaptation.

\subsection{Scope}

This evaluation focuses on I/O-level data manipulation---the attacker reads/writes PLC I/O tags via standard S7 protocol. Written values to PLC-refreshed tags are overwritten each scan cycle, requiring precise timing or continuous injection. The safety shield bounds all writes to process-safe ranges, so reported attack success rates reflect effectiveness within this configured envelope. The framework does not model firmware compromise, PLC root access, or unrestricted code deployment.
